% Part: first-order-logic
% Chapter: natural-deduction
% Section: provability-quantifiers

% verification of properties of provability needed for maximally
% consistent sets in the completeness chapter.

\documentclass[../../../include/open-logic-section]{subfiles}

\begin{document}

\olfileid[id]{fol}{ntd}{qpr}

\olsection{\usetoken{S}{derivability} dan Kuantor}

\begin{explain}
  Teorema kelengkapan juga memerlukan agar aturan-aturan deduksi alami
  menghasilkan fakta-fakta tentang~$\Proves$ yang dibuktikan dalam
  bagian ini.
\end{explain}

\begin{thm}
\ollabel{thm:strong-generalization} Jika $c$ adalah simbol konstanta yang
tidak muncul dalam $\Gamma$ ataupun $!A(x)$ dan $\Gamma \Proves !A(c)$,
maka $\Gamma \Proves \lforall[x][!A(x)]$.
\end{thm}

\begin{proof}
Misalkan $\delta$ adalah !!a{derivation} untuk $!A(c)$ dari $\Gamma$.
Dengan menambahkan inferensi \Intro{\lforall}, kita memperoleh
!!a{derivation} untuk $\lforall[x][!A(x)]$. Karena $c$ tidak muncul dalam
$\Gamma$ ataupun $!A(x)$, syarat variabel eigen terpenuhi.
\end{proof}

\begin{prop}
\ollabel{prop:provability-quantifiers}
\begin{tagenumerate}{prvEx,prvAll}
\tagitem{prvEx}{$!A(t) \Proves \lexists[x][!A(x)]$.}{}

\tagitem{prvAll}{$\lforall[x][!A(x)] \Proves !A(t)$.}{}
\end{tagenumerate}
\end{prop}

\begin{proof}
\begin{tagenumerate}{prvEx,prvAll}
\tagitem{prvEx}{Berikut ini merupakan !!a{derivation} untuk
  $\lexists[x][!A(x)]$ dari~$!A(t)$:
\begin{prooftree}
\AxiomC{$!A(t)$}
\RightLabel{\Intro{\lexists}}
\UnaryInfC{$\lexists[x][!A(x)]$}
\end{prooftree}}{}

\tagitem{prvAll}{Berikut ini merupakan !!a{derivation} untuk~$!A(t)$
  dari~$\lforall[x][!A(x)]$:
\begin{prooftree}
\AxiomC{$\lforall[x][!A(x)]$}
\RightLabel{\Elim{\lforall}}
\UnaryInfC{$!A(t)$}
\end{prooftree}}{}
\end{tagenumerate}
\end{proof}

\end{document}
